
\documentclass{article}
\usepackage{colortbl}
\usepackage{makecell}
\usepackage{multirow}
\usepackage{supertabular}

\begin{document}

\newcounter{utterance}

\centering \large Interaction Transcript for game `function\_detective', experiment `string\_passive\_labeled\_pairs', episode 171 with Qwen3.6{-}35B{-}A3B{-}FP8{-}without{-}reasoning.
\vspace{24pt}

{ \footnotesize  \setcounter{utterance}{1}
\setlength{\tabcolsep}{0pt}
\begin{supertabular}{c@{$\;$}|p{.15\linewidth}@{}p{.15\linewidth}p{.15\linewidth}p{.15\linewidth}p{.15\linewidth}p{.15\linewidth}}
    \# & \multicolumn{2}{c}{Player} && \multicolumn{2}{c}{Game Master} \\
    \hline

    \theutterance \stepcounter{utterance}  
    & & & \multicolumn{4}{p{0.6\linewidth}}{
        \cellcolor[rgb]{0.9,0.9,0.9}{
            \makecell[{{p{\linewidth}}}]{
                \texttt{\tiny{[P1$\langle$GM]}}
                \texttt{I have a hidden Python function for you to figure out.} \\
\\ 
\texttt{Your job is to infer what it does and then submit a correct Python implementation.} \\
\\ 
\texttt{Function signature:} \\
\texttt{(text: str) {-}> str} \\
\\ 
\texttt{You have 15 turns total.} \\
\\ 
\texttt{I reveal candidate input{-}output pairs one at a time and tell you whether each one is correct.} \\
\\ 
\texttt{You must make your best SOLVE attempt no later than the second{-}to{-}last turn.} \\
\texttt{Do not wait until the very end.} \\
\\ 
\texttt{You can use these actions:} \\
\texttt{{-} JSON format described at the end of the prompt} \\
\texttt{{-} SOLVE: ```python ... ```} \\
\\ 
\texttt{Do not add import statements in your guess.} \\
\texttt{Anything you might need is already available, so just write the function body in hypothesis and return the answer.} \\
\\ 
\texttt{Preloaded examples:} \\
\texttt{1. ((' ', ''), True)} \\
\\ 
\texttt{Scoring:} \\
\texttt{{-} correctness is required for any non{-}zero score} \\
\texttt{{-} in interactive modes, solving with fewer information requests is better} \\
\\ 
\texttt{Input formatting reminders:} \\
\texttt{{-} strings must be quoted} \\
\texttt{{-} booleans must be True/False} \\
\texttt{{-} use commas between multiple values} \\
\texttt{{-} do not write any extra action tags in the same message} \\
\\ 
\texttt{Format example:} \\
\texttt{You write:} \\
\texttt{\{'current\_mode': <probing/confirming/don't know; choose whichever best describes your current state in the task>, 'input\_rationale': <brief description of why you chose this hypothesis>, 'hypothesis': <either the exact string "unknown" if you have no guess yet OR your current best guess as runnable Python matching the given signature, with a one{-}line docstring stating what it does. Never write placeholder code, if you cannot commit to a guess, use "unknown". Match the given signature.} \\
\texttt{e.g.} \\
\texttt{def f(x):} \\
\texttt{	"""Returns the absolute value of x."""} \\
\texttt{	return x if x >= 0 else {-}x} \\
\\ 
\texttt{Use "unknown" if you have no guess.>, 'confidence\_score': <score between 1{-}100. 100 means highest confidence in your hypothesis>\}} \\
\\ 
\texttt{Output only one valid json per turn.} \\
            }
        }
    }
     \\ \\

    \theutterance \stepcounter{utterance}  
    & \multicolumn{4}{p{0.6\linewidth}}{
        \cellcolor[rgb]{0.561,0.737,0.561}{
            \makecell[{{p{\linewidth}}}]{
                \texttt{\tiny{[P1$\rangle$GM]}}
                \texttt{\{'current\_mode': 'probing', 'input\_rationale': 'I have only one example so far: a single space maps to an empty string. This is very sparse. I need more data points to understand the transformation. I will try a few more inputs: an empty string, a single letter, and a string with multiple spaces to see if it strips whitespace, removes specific characters, or does something else.', 'hypothesis': 'unknown', 'confidence\_score': 10\}} \\
            }
        }
    }
    & & \\ \\

    \theutterance \stepcounter{utterance}  
    & & & \multicolumn{4}{p{0.6\linewidth}}{
        \cellcolor[rgb]{0.9,0.9,0.9}{
            \makecell[{{p{\linewidth}}}]{
                \texttt{\tiny{[P1$\langle$GM]}}
                \texttt{Game Violation: ParseError: Response is not valid JSON.} \\
\texttt{Output ONLY:} \\
\texttt{One JSON with confidence\_score, current\_mode, input\_rationale, and hypothesis keys ...} \\
\texttt{OR} \\
\texttt{SOLVE: ```python ... ```} \\
            }
        }
    }
     \\ \\

    \theutterance \stepcounter{utterance}  
    & \multicolumn{4}{p{0.6\linewidth}}{
        \cellcolor[rgb]{0.561,0.737,0.561}{
            \makecell[{{p{\linewidth}}}]{
                \texttt{\tiny{[P1$\rangle$GM]}}
                \texttt{\{"current\_mode": "probing", "input\_rationale": "I have only one example: (' ', '') {-}> True. This is extremely sparse. I need to test more inputs to understand the pattern. I will test an empty string, a single non{-}space character, and a string with multiple spaces.", "hypothesis": "unknown", "confidence\_score": 10\}} \\
            }
        }
    }
    & & \\ \\

    \theutterance \stepcounter{utterance}  
    & & & \multicolumn{4}{p{0.6\linewidth}}{
        \cellcolor[rgb]{0.9,0.9,0.9}{
            \makecell[{{p{\linewidth}}}]{
                \texttt{\tiny{[P1$\langle$GM]}}
                \texttt{OUTPUT: (('   ', '\_alt'), False)} \\
            }
        }
    }
     \\ \\

    \theutterance \stepcounter{utterance}  
    & \multicolumn{4}{p{0.6\linewidth}}{
        \cellcolor[rgb]{0.561,0.737,0.561}{
            \makecell[{{p{\linewidth}}}]{
                \texttt{\tiny{[P1$\rangle$GM]}}
                \texttt{\{"current\_mode": "probing", "input\_rationale": "The first example was (' ', '') {-}> True. The second example is ('   ', '\_alt') {-}> False. This suggests the function might be checking if the input string is a single space character. If it is, it returns an empty string (or True in the context of the test framework, perhaps indicating the output matches an expected empty string). If it's not a single space, it returns something else (or False). Let's test with a single non{-}space character to see if it returns that character or something else.", "hypothesis": "unknown", "confidence\_score": 20\}} \\
            }
        }
    }
    & & \\ \\

    \theutterance \stepcounter{utterance}  
    & & & \multicolumn{4}{p{0.6\linewidth}}{
        \cellcolor[rgb]{0.9,0.9,0.9}{
            \makecell[{{p{\linewidth}}}]{
                \texttt{\tiny{[P1$\langle$GM]}}
                \texttt{OUTPUT: (('a', 'a'), True)} \\
            }
        }
    }
     \\ \\

    \theutterance \stepcounter{utterance}  
    & \multicolumn{4}{p{0.6\linewidth}}{
        \cellcolor[rgb]{0.561,0.737,0.561}{
            \makecell[{{p{\linewidth}}}]{
                \texttt{\tiny{[P1$\rangle$GM]}}
                \texttt{\{"current\_mode": "probing", "input\_rationale": "So far: (' ', '') {-}> True, ('   ', '\_alt') {-}> False, ('a', 'a') {-}> True. It seems like for single characters, it returns the character itself. For a single space, it returns an empty string. For multiple spaces, it returns '\_alt'. This is confusing. Let's test more cases: a two{-}letter string, a string with a space and a letter, and an empty string.", "hypothesis": "unknown", "confidence\_score": 30\}} \\
            }
        }
    }
    & & \\ \\

    \theutterance \stepcounter{utterance}  
    & & & \multicolumn{4}{p{0.6\linewidth}}{
        \cellcolor[rgb]{0.9,0.9,0.9}{
            \makecell[{{p{\linewidth}}}]{
                \texttt{\tiny{[P1$\langle$GM]}}
                \texttt{OUTPUT: (('hello world', 'helloworld\_alt'), False)} \\
            }
        }
    }
     \\ \\

    \theutterance \stepcounter{utterance}  
    & \multicolumn{4}{p{0.6\linewidth}}{
        \cellcolor[rgb]{0.561,0.737,0.561}{
            \makecell[{{p{\linewidth}}}]{
                \texttt{\tiny{[P1$\rangle$GM]}}
                \texttt{\{"current\_mode": "probing", "input\_rationale": "Let's analyze the examples:\textbackslash{}n1. (' ', '') {-}> True: Input is a single space. Output is empty string. The test says True, meaning the actual output of the hidden function for ' ' is ''. \textbackslash{}n2. ('   ', '\_alt') {-}> False: Input is three spaces. Output is '\_alt'. The test says False, meaning the actual output of the hidden function for '   ' is NOT '\_alt'.\textbackslash{}n3. ('a', 'a') {-}> True: Input is 'a'. Output is 'a'. The test says True, meaning the actual output of the hidden function for 'a' is 'a'.\textbackslash{}n4. ('hello world', 'helloworld\_alt') {-}> False: Input is 'hello world'. Output is 'helloworld\_alt'. The test says False, meaning the actual output of the hidden function for 'hello world' is NOT 'helloworld\_alt'.\textbackslash{}n\textbackslash{}nFrom 1 and 3, it seems like the function might be removing spaces. ' ' {-}> '' (removes the space). 'a' {-}> 'a' (no space to remove). \textbackslash{}nLet's check if 'hello world' {-}> 'helloworld'. If so, then the test ('hello world', 'helloworld\_alt') {-}> False makes sense because 'helloworld' != 'helloworld\_alt'.\textbackslash{}nWhat about '   ' {-}> ''? If} \\
            }
        }
    }
    & & \\ \\

    \theutterance \stepcounter{utterance}  
    & & & \multicolumn{4}{p{0.6\linewidth}}{
        \cellcolor[rgb]{0.9,0.9,0.9}{
            \makecell[{{p{\linewidth}}}]{
                \texttt{\tiny{[P1$\langle$GM]}}
                \texttt{Game Violation: ParseError: Response is not valid JSON.} \\
\texttt{Output ONLY:} \\
\texttt{One JSON with confidence\_score, current\_mode, input\_rationale, and hypothesis keys ...} \\
\texttt{OR} \\
\texttt{SOLVE: ```python ... ```} \\
            }
        }
    }
     \\ \\

    \theutterance \stepcounter{utterance}  
    & \multicolumn{4}{p{0.6\linewidth}}{
        \cellcolor[rgb]{0.561,0.737,0.561}{
            \makecell[{{p{\linewidth}}}]{
                \texttt{\tiny{[P1$\rangle$GM]}}
                \texttt{\{"current\_mode": "probing", "input\_rationale": "Based on previous outputs: (' ', '') {-}> True implies f(' ') == ''. ('a', 'a') {-}> True implies f('a') == 'a'. ('   ', '\_alt') {-}> False implies f('   ') != '\_alt'. ('hello world', 'helloworld\_alt') {-}> False implies f('hello world') != 'helloworld\_alt'. The pattern suggests f(s) removes all spaces from s. Let's verify this hypothesis with a new test case that would be True if f(s) == s.replace(' ', ''). For example, ('ab cd', 'abcd') should be True. If f(s) == s.replace(' ', ''), then f('ab cd') is 'abcd'.", "hypothesis": "unknown", "confidence\_score": 60\}} \\
            }
        }
    }
    & & \\ \\

    \theutterance \stepcounter{utterance}  
    & & & \multicolumn{4}{p{0.6\linewidth}}{
        \cellcolor[rgb]{0.9,0.9,0.9}{
            \makecell[{{p{\linewidth}}}]{
                \texttt{\tiny{[P1$\langle$GM]}}
                \texttt{OUTPUT: (('123', ''), True)} \\
            }
        }
    }
     \\ \\

    \theutterance \stepcounter{utterance}  
    & \multicolumn{4}{p{0.6\linewidth}}{
        \cellcolor[rgb]{0.561,0.737,0.561}{
            \makecell[{{p{\linewidth}}}]{
                \texttt{\tiny{[P1$\rangle$GM]}}
                \texttt{\{"current\_mode": "probing", "input\_rationale": "The previous hypothesis that f(s) == s.replace(' ', '') is incorrect because f('123') is '' according to the new output, but '123'.replace(' ', '') is '123'. So the function does something else. Let's re{-}evaluate all examples:\textbackslash{}n1. f(' ') == ''\textbackslash{}n2. f('   ') != '\_alt' (likely f('   ') == '' or something else)\textbackslash{}n3. f('a') == 'a'\textbackslash{}n4. f('hello world') != 'helloworld\_alt' (likely f('hello world') == 'helloworld' or something else)\textbackslash{}n5. f('123') == ''\textbackslash{}n\textbackslash{}nThis is very strange. 'a' stays 'a', but '123' becomes ''. ' ' becomes ''. 'hello world' likely becomes 'helloworld' (if it removes spaces). But '123' has no spaces and becomes empty. 'a' has no spaces and stays 'a'.\textbackslash{}n\textbackslash{}nMaybe it depends on the character type? Letters stay, digits disappear? Spaces disappear?\textbackslash{}nLet's test with a mix: 'a1' and '1a'. Also test 'ab' and '12'.", "hypothesis": "unknown", "confidence\_score": 40\}} \\
            }
        }
    }
    & & \\ \\

    \theutterance \stepcounter{utterance}  
    & & & \multicolumn{4}{p{0.6\linewidth}}{
        \cellcolor[rgb]{0.9,0.9,0.9}{
            \makecell[{{p{\linewidth}}}]{
                \texttt{\tiny{[P1$\langle$GM]}}
                \texttt{OUTPUT: (('ab', 'ab\_alt'), False)} \\
            }
        }
    }
     \\ \\

    \theutterance \stepcounter{utterance}  
    & \multicolumn{4}{p{0.6\linewidth}}{
        \cellcolor[rgb]{0.561,0.737,0.561}{
            \makecell[{{p{\linewidth}}}]{
                \texttt{\tiny{[P1$\rangle$GM]}}
                \texttt{\{"current\_mode": "probing", "input\_rationale": "Let's re{-}analyze all confirmed outputs:\textbackslash{}n1. f(' ') == ''\textbackslash{}n2. f('a') == 'a'\textbackslash{}n3. f('123') == ''\textbackslash{}n4. f('ab') != 'ab\_alt' (likely f('ab') == 'ab' based on pattern with 'a')\textbackslash{}n\textbackslash{}nHypothesis: The function keeps alphabetic characters and removes non{-}alphabetic characters (spaces, digits). \textbackslash{}n{-} ' ' (space) {-}> removed {-}> ''\textbackslash{}n{-} 'a' (alpha) {-}> kept {-}> 'a'\textbackslash{}n{-} '123' (digits) {-}> removed {-}> ''\textbackslash{}n{-} 'ab' (alpha) {-}> kept {-}> 'ab'\textbackslash{}n{-} 'hello world' {-}> 'helloworld' (spaces removed, letters kept)\textbackslash{}n\textbackslash{}nLet's test this with a string containing both letters and digits, e.g., 'a1b'. If the hypothesis is correct, f('a1b') should be 'ab'.", "hypothesis": "unknown", "confidence\_score": 70\}} \\
            }
        }
    }
    & & \\ \\

    \theutterance \stepcounter{utterance}  
    & & & \multicolumn{4}{p{0.6\linewidth}}{
        \cellcolor[rgb]{0.9,0.9,0.9}{
            \makecell[{{p{\linewidth}}}]{
                \texttt{\tiny{[P1$\langle$GM]}}
                \texttt{OUTPUT: (('a{-}b', 'ab'), True)} \\
            }
        }
    }
     \\ \\

    \theutterance \stepcounter{utterance}  
    & \multicolumn{4}{p{0.6\linewidth}}{
        \cellcolor[rgb]{0.561,0.737,0.561}{
            \makecell[{{p{\linewidth}}}]{
                \texttt{\tiny{[P1$\rangle$GM]}}
                \texttt{\{"current\_mode": "probing", "input\_rationale": "New data point: f('a{-}b') == 'ab'. This supports the hypothesis that non{-}alphabetic characters are removed. 'a' and 'b' are alphabetic, so they are kept. '{-}' is not alphabetic, so it is removed. Result is 'ab'.\textbackslash{}n\textbackslash{}nLet's verify the previous hypotheses:\textbackslash{}n{-} f(' ') == '' (space is not alpha, removed)\textbackslash{}n{-} f('a') == 'a' (alpha, kept)\textbackslash{}n{-} f('123') == '' (digits not alpha, removed)\textbackslash{}n{-} f('ab') == 'ab' (alphas kept, so f('ab') != 'ab\_alt' is True because 'ab' != 'ab\_alt')\textbackslash{}n{-} f('hello world') == 'helloworld' (spaces removed, letters kept, so f('hello world') != 'helloworld\_alt' is True because 'helloworld' != 'helloworld\_alt')\textbackslash{}n\textbackslash{}nThe hypothesis seems robust: The function returns the input string with all non{-}alphabetic characters removed.\textbackslash{}n\textbackslash{}nI will test one more case to be sure: a string with only non{-}alphabetic characters, e.g., '1{-}2', which should result in ''. And a string with mixed, e.g., 'a1b2c', which should result in 'abc'.", "hypothesis": "unknown",} \\
            }
        }
    }
    & & \\ \\

    \theutterance \stepcounter{utterance}  
    & & & \multicolumn{4}{p{0.6\linewidth}}{
        \cellcolor[rgb]{0.9,0.9,0.9}{
            \makecell[{{p{\linewidth}}}]{
                \texttt{\tiny{[P1$\langle$GM]}}
                \texttt{Game Violation: ParseError: Response is not valid JSON.} \\
\texttt{Output ONLY:} \\
\texttt{One JSON with confidence\_score, current\_mode, input\_rationale, and hypothesis keys ...} \\
\texttt{OR} \\
\texttt{SOLVE: ```python ... ```} \\
            }
        }
    }
     \\ \\

    \theutterance \stepcounter{utterance}  
    & \multicolumn{4}{p{0.6\linewidth}}{
        \cellcolor[rgb]{0.561,0.737,0.561}{
            \makecell[{{p{\linewidth}}}]{
                \texttt{\tiny{[P1$\rangle$GM]}}
                \texttt{\{"current\_mode": "confirming", "input\_rationale": "The hypothesis is that the function removes all non{-}alphabetic characters from the string. This fits all observed data: ' ' {-}> '', 'a' {-}> 'a', '123' {-}> '', 'ab' {-}> 'ab', 'a{-}b' {-}> 'ab'. I will test '1{-}2' which should result in '' to confirm non{-}alpha removal, and 'a1b2c' which should result in 'abc'.", "hypothesis": "unknown", "confidence\_score": 90\}} \\
            }
        }
    }
    & & \\ \\

    \theutterance \stepcounter{utterance}  
    & & & \multicolumn{4}{p{0.6\linewidth}}{
        \cellcolor[rgb]{0.9,0.9,0.9}{
            \makecell[{{p{\linewidth}}}]{
                \texttt{\tiny{[P1$\langle$GM]}}
                \texttt{OUTPUT: (('hello', 'hello\_alt'), False)} \\
            }
        }
    }
     \\ \\

    \theutterance \stepcounter{utterance}  
    & \multicolumn{4}{p{0.6\linewidth}}{
        \cellcolor[rgb]{0.561,0.737,0.561}{
            \makecell[{{p{\linewidth}}}]{
                \texttt{\tiny{[P1$\rangle$GM]}}
                \texttt{\{"current\_mode": "confirming", "input\_rationale": "The output ('hello', 'hello\_alt') {-}> False confirms that f('hello') is not 'hello\_alt'. Since 'hello' consists only of alphabetic characters, and our hypothesis is to keep only alphabetic characters, f('hello') should be 'hello'. Since 'hello' != 'hello\_alt', the test returns False, which is consistent. This further strengthens the hypothesis that the function filters out non{-}alphabetic characters. I am now confident enough to submit the solution.", "hypothesis": "def f(text: str) {-}> str:\textbackslash{}n    \textbackslash{}"\textbackslash{}"\textbackslash{}"Returns the input string with all non{-}alphabetic characters removed.\textbackslash{}"\textbackslash{}"\textbackslash{}"\textbackslash{}n    return ''.join(char for char in text if char.isalpha())", "confidence\_score": 95\}} \\
            }
        }
    }
    & & \\ \\

    \theutterance \stepcounter{utterance}  
    & & & \multicolumn{4}{p{0.6\linewidth}}{
        \cellcolor[rgb]{0.9,0.9,0.9}{
            \makecell[{{p{\linewidth}}}]{
                \texttt{\tiny{[P1$\langle$GM]}}
                \texttt{OUTPUT: (('aA', 'aA'), True)} \\
            }
        }
    }
     \\ \\

    \theutterance \stepcounter{utterance}  
    & \multicolumn{4}{p{0.6\linewidth}}{
        \cellcolor[rgb]{0.561,0.737,0.561}{
            \makecell[{{p{\linewidth}}}]{
                \texttt{\tiny{[P1$\rangle$GM]}}
                \texttt{\{"current\_mode": "confirming", "input\_rationale": "The output ('aA', 'aA') {-}> True confirms that f('aA') == 'aA'. Since 'a' and 'A' are both alphabetic, they are kept. This is consistent with the hypothesis that the function removes all non{-}alphabetic characters. I have enough evidence to submit the solution.", "hypothesis": "def f(text: str) {-}> str:\textbackslash{}n    \textbackslash{}"\textbackslash{}"\textbackslash{}"Returns the input string with all non{-}alphabetic characters removed.\textbackslash{}"\textbackslash{}"\textbackslash{}"\textbackslash{}n    return ''.join(char for char in text if char.isalpha())", "confidence\_score": 98\}} \\
            }
        }
    }
    & & \\ \\

    \theutterance \stepcounter{utterance}  
    & & & \multicolumn{4}{p{0.6\linewidth}}{
        \cellcolor[rgb]{0.9,0.9,0.9}{
            \makecell[{{p{\linewidth}}}]{
                \texttt{\tiny{[P1$\langle$GM]}}
                \texttt{OUTPUT: (('', '\_alt'), False)} \\
            }
        }
    }
     \\ \\

    \theutterance \stepcounter{utterance}  
    & \multicolumn{4}{p{0.6\linewidth}}{
        \cellcolor[rgb]{0.561,0.737,0.561}{
            \makecell[{{p{\linewidth}}}]{
                \texttt{\tiny{[P1$\rangle$GM]}}
                \texttt{\{"current\_mode": "confirming", "input\_rationale": "The output ('', '\_alt') {-}> False confirms that f('') is not '\_alt'. Since the empty string has no characters, f('') should be ''. Since '' != '\_alt', the test returns False. This is consistent with the hypothesis. I am now fully confident in the solution.", "hypothesis": "def f(text: str) {-}> str:\textbackslash{}n    \textbackslash{}"\textbackslash{}"\textbackslash{}"Returns the input string with all non{-}alphabetic characters removed.\textbackslash{}"\textbackslash{}"\textbackslash{}"\textbackslash{}n    return ''.join(char for char in text if char.isalpha())", "confidence\_score": 100\}} \\
            }
        }
    }
    & & \\ \\

    \theutterance \stepcounter{utterance}  
    & & & \multicolumn{4}{p{0.6\linewidth}}{
        \cellcolor[rgb]{0.9,0.9,0.9}{
            \makecell[{{p{\linewidth}}}]{
                \texttt{\tiny{[P1$\langle$GM]}}
                \texttt{OUTPUT: (('A', 'A'), True)} \\
            }
        }
    }
     \\ \\

    \theutterance \stepcounter{utterance}  
    & \multicolumn{4}{p{0.6\linewidth}}{
        \cellcolor[rgb]{0.561,0.737,0.561}{
            \makecell[{{p{\linewidth}}}]{
                \texttt{\tiny{[P1$\rangle$GM]}}
                \texttt{\{"current\_mode": "confirming", "input\_rationale": "The output ('A', 'A') {-}> True confirms that f('A') == 'A'. This is consistent with the hypothesis that the function keeps only alphabetic characters. All previous tests also support this hypothesis. I am now fully confident and ready to submit the final solution.", "hypothesis": "def f(text: str) {-}> str:\textbackslash{}n    \textbackslash{}"\textbackslash{}"\textbackslash{}"Returns the input string with all non{-}alphabetic characters removed.\textbackslash{}"\textbackslash{}"\textbackslash{}"\textbackslash{}n    return ''.join(char for char in text if char.isalpha())", "confidence\_score": 100\}} \\
            }
        }
    }
    & & \\ \\

    \theutterance \stepcounter{utterance}  
    & & & \multicolumn{4}{p{0.6\linewidth}}{
        \cellcolor[rgb]{0.9,0.9,0.9}{
            \makecell[{{p{\linewidth}}}]{
                \texttt{\tiny{[P1$\langle$GM]}}
                \texttt{OUTPUT: (('aba', 'aba\_alt'), False)} \\
            }
        }
    }
     \\ \\

    \theutterance \stepcounter{utterance}  
    & \multicolumn{4}{p{0.6\linewidth}}{
        \cellcolor[rgb]{0.561,0.737,0.561}{
            \makecell[{{p{\linewidth}}}]{
                \texttt{\tiny{[P1$\rangle$GM]}}
                \texttt{\{"current\_mode": "confirming", "input\_rationale": "The output ('aba', 'aba\_alt') {-}> False confirms that f('aba') is not 'aba\_alt'. Since 'aba' consists only of alphabetic characters, f('aba') should be 'aba'. Since 'aba' != 'aba\_alt', the test returns False. This is consistent with the hypothesis. I have enough evidence to submit the solution.", "hypothesis": "def f(text: str) {-}> str:\textbackslash{}n    \textbackslash{}"\textbackslash{}"\textbackslash{}"Returns the input string with all non{-}alphabetic characters removed.\textbackslash{}"\textbackslash{}"\textbackslash{}"\textbackslash{}n    return ''.join(char for char in text if char.isalpha())", "confidence\_score": 100\}} \\
            }
        }
    }
    & & \\ \\

    \theutterance \stepcounter{utterance}  
    & & & \multicolumn{2}{p{0.3\linewidth}}{
        \cellcolor[rgb]{0.9,0.9,0.9}{
            \makecell[{{p{\linewidth}}}]{
                \texttt{\tiny{[GM$|$GM]}}
                \texttt{Game Verdict: LOSS (turn limit reached)} \\
            }
        }
    }
    & & \\ \\

\end{supertabular}
}

\end{document}
